\documentclass[conference]{IEEEtran}
\IEEEoverridecommandlockouts
\usepackage{cite}
\usepackage{amsmath,amssymb,amsfonts}
\usepackage{graphicx}
\usepackage{booktabs}
\usepackage{multirow}
\usepackage{xcolor}
\usepackage{url}
\usepackage{stfloats}
\newcommand{\best}[1]{\textbf{#1}}
\newcommand{\ours}{TypeAnchor}
\newcommand{\fig}[2]{\fbox{\parbox[c][#1][c]{\dimexpr\linewidth-2\fboxsep-2\fboxrule\relax}{\centering\small #2}}}

\begin{document}

\title{\ours{}: Text-Anchored Multi-Type Defect Segmentation on a Frozen Vision-Language Model for Concrete Bridge Inspection}

\author{\IEEEauthorblockN{Qiang Liu$^{1,2}$, Haifeng Zhong$^{1,2}$}
\IEEEauthorblockA{$^{1}$School of Computer Science and Artificial Intelligence, Hunan University of Technology, Zhuzhou, China\\
$^{2}$Hunan Provincial Engineering Technology Research Center of Industrial Data Intelligence, Zhuzhou, China\\
liuqiang@hut.edu.cn, hfz422875@stu.hut.edu.cn}}

\maketitle

\begin{abstract}
Repair decisions for concrete bridges depend on the defect type: cracks are measured by width and orientation, spalling and corrosion by exposed area, efflorescence by the moisture path behind it. A pixel map that only separates damaged from intact surface is not enough, and the types overlap on the same surface. Fully supervised segmentation networks deliver type-specific masks but lose accuracy on images from another agency, while anomaly detectors on frozen CLIP transfer better but return one binary map. We present \ours{}, a light supervised head that uses the text embeddings of a frozen CLIP ViT-L/14 as a class interface. Each defect type receives a text anchor; an Anchor Margin Readout turns the similarity margin between the type anchor and the intact-surface anchor into an independent per-type probability, so types may overlap. A Local-Regional Similarity Mixture learns per anchor how much patch-level and neighbourhood similarity to use, which keeps thin cracks visible while large defects draw on context. A Reference-Residual Gate derived from one intact image scales a class-agnostic defect channel and the image score. With 6.6 M trainable parameters, \ours{} reaches 76.8 binary pixel AP on the held-out CODEBRIM test split, level with a fully trained SegFormer-B2 and above WPFormer and four CLIP-based detectors trained on the same tiles, and at 90\% pixel recall it flags less than half the intact-surface area of every other method. On S2DS and dacl10k without target data, type-aware frozen-CLIP heads exceed both supervised networks, and the simpler MultiADS exceeds \ours{}: in-domain accuracy of a supervised head is accuracy on the training domain, while transfer comes from the frozen backbone and its text alignment.
\end{abstract}

\begin{IEEEkeywords}
bridge inspection, defect segmentation, vision-language models, CLIP, multi-label segmentation, cross-dataset evaluation
\end{IEEEkeywords}

\section{Introduction}
Visual inspection of concrete bridges produces images in which several kinds of damage coexist: cracks as narrow discontinuities, spalling as broken-off regions that expose aggregate or reinforcement, corrosion as rust staining around exposed steel, and efflorescence as white deposits left by water moving through the concrete. Maintenance acts on these types differently. A crack is assessed by its width and orientation, spalling and corrosion by the exposed area and the condition of the reinforcement, and efflorescence as evidence of a leakage path. The pixel map that an inspection system hands to an engineer therefore needs to say which type is where, and it needs to allow overlap, because corrosion appears inside spalled regions and efflorescence runs across cracks. The map also needs to survive a change of image collection. Agencies photograph different structures with different cameras and lighting, and pixel-level annotation of every new collection is the most expensive step in deploying such a system.

Two families of methods meet parts of this requirement. Fully supervised segmentation networks, from general architectures such as SegFormer~\cite{segformer} to surface-defect designs such as WPFormer~\cite{wpformer} and crack-specific networks~\cite{scsegamba}, give type-specific masks and high accuracy on the collection they were trained on. On a different collection their accuracy drops, and they are silent about damage outside the training classes. Anomaly detection on frozen vision-language models moves in the opposite direction. AnomalyCLIP~\cite{anomalyclip}, AdaptCLIP~\cite{adaptclip} and Bayes-PFL~\cite{bayespfl} keep the CLIP~\cite{clip} encoders fixed, learn a few million parameters that align patch features with normal and abnormal text, and transfer across industrial domains with zero or one normal image. Their output is a single binary anomaly map. MultiADS~\cite{multiads} extends this line with one text prompt per defect type and a pixel-level softmax over the types, which produces type maps but assigns each pixel one type and was designed on industrial objects photographed under controlled conditions.

Three obstacles separate these families from the inspection requirement. The class interface is the first. A normal-versus-abnormal formulation has no place for types, and a softmax over types has no place for overlap. The second is scale. At the 518-pixel input of CLIP ViT-L/14 the patch grid is $37\times37$, about 14 pixels per cell, and cracks are often thinner than one cell. Pooling similarity over a neighbourhood, which stabilises large defects, erases cracks: in our ablation, replacing patch-level similarity by $5\times5$ regional pooling drops crack AP from 35.7 to 8.2 while spalling and efflorescence are unaffected. Scale requirements differ by type. The third is measurement. Public civil-engineering datasets store flattened single-label masks in which a priority rule removes overlap, sets labelled as intact contain unlabelled defects that inflate false-alarm measurements, two common implementations of P-AUPRO disagree by 4.8 points on identical predictions, and the defect-size distribution of the evaluation split moves the margin between two methods by more than ten points.

\ours{} addresses the interface and scale obstacles with one design decision: the text anchors of a frozen CLIP serve as a class interface for a light supervised head, not as an anomaly detector. Six anchors represent intact surface, the four defect types and generic damaged surface. The Local-Regional Similarity Mixture (LRM) computes, for each anchor, the similarity to the projected patch tokens at patch level and over a residual-weighted neighbourhood and mixes the two with a learned per-anchor weight, so the crack anchor can stay local while the spalling and efflorescence anchors use context. The Anchor Margin Readout (AMR) passes the margin between each type anchor and the intact anchor through an independent sigmoid; a pixel may be positive for several types. The Reference-Residual Gate (RRG) keeps the residual between the query tokens and the nearest tokens of one intact reference image as a gate that scales a class-agnostic defect channel and the image-level head. The head has 6.6 M trainable parameters and needs one arbitrary intact image at inference.

We trained \ours{} on 2{,}861 CODEBRIM~\cite{codebrim} tiles, 917 of which carry overlapping per-type masks that we annotated, selected every design choice on the validation split, and report on a test split of 809 tiles that no choice touched, and, without target data, on S2DS~\cite{s2ds} and a four-type subset of dacl10k~\cite{dacl10k}. On the test split the head reaches 76.8 binary pixel AP over three seeds, level with SegFormer-B2 at 76.4 and above WPFormer, AdaptCLIP, MultiADS, Bayes-PFL and AnomalyCLIP trained with their official code on the same tiles. At 90\% pixel recall it marks 4.4\% of the intact surface, against 9.5\% for AdaptCLIP and 10.0\% for SegFormer-B2, and its per-type mean AP of 59.4 compares with 56.2 for MultiADS and 60.1 for SegFormer-B2. Across datasets the picture inverts: \ours{} and MultiADS exceed both supervised networks in pixel AP on S2DS and dacl10k, and MultiADS, with 3.1 M linear parameters, exceeds \ours{}. The evaluation split matters as much as the method: on the validation split, where the median defect covers 6{,}300 pixels, \ours{} is 20 points of pixel AP ahead of AdaptCLIP; on the test split, where it covers 44{,}000, the margin is 3.5.

Our contributions are as follows. We propose \ours{}, a text-anchored multi-type segmentation head for frozen CLIP with independent per-type readouts on anchor margins, a per-anchor local-regional similarity mixture and a residual-gated class-agnostic channel, and we show through ablation which obstacle each part removes. We provide overlapping per-type pixel masks for CODEBRIM tiles together with an evaluation protocol for bridge inspection: per-type multi-label metrics, false-alarm area at matched recall, a cleaned validation split, two external test sets and a reconciled P-AUPRO. We report, on one protocol, how CLIP-based detectors, their multi-type extension and two fully supervised networks compare in-domain and across datasets, and trace in-domain accuracy and cross-dataset accuracy to different causes.

\section{Related Work}
\textbf{Defect segmentation in civil structures.} Pixel-level bridge inspection has been served by general segmentation backbones trained on inspection data~\cite{segformer,dacl10k}, by query-based decoders for surface defects~\cite{wpformer}, and by lightweight crack networks that exploit the thin, connected structure of cracks~\cite{scsegamba}. These networks learn the appearance of the training collection and return one label per pixel. CODEBRIM~\cite{codebrim} introduced multi-target labels at image level for six defect types; S2DS~\cite{s2ds} and dacl10k~\cite{dacl10k} provide pixel masks with a priority rule that flattens overlapping damage into one class per pixel. We annotate overlapping per-type masks on CODEBRIM tiles and evaluate each type as its own binary problem.

\textbf{Anomaly detection on frozen CLIP.} WinCLIP~\cite{winclip} scored patches against hand-written normal and abnormal prompts; AnomalyCLIP~\cite{anomalyclip} learned object-agnostic prompts; AdaptCLIP~\cite{adaptclip} added visual and textual adapters and a prompt-query residual branch that compares the query with one normal image; Bayes-PFL~\cite{bayespfl} samples prompts from a learned distribution; AA-CLIP~\cite{aaclip} sharpens the normal-abnormal separation of the text space. All of these output a binary anomaly map and an image score. MultiADS~\cite{multiads} is the closest method to ours: it keeps the backbone frozen, trains linear projections, and scores pixels with a softmax over a normal prompt and several defect-type prompts. \ours{} differs in the readout, which is a set of independent sigmoids on anchor margins, in the per-anchor local-regional mixture, and in the residual gate; Section~\ref{sec:exp} compares the two under one protocol. P-AUPRO follows the per-region overlap of~\cite{bergmann2020}; we use the GPU implementation of anomalib~\cite{anomalib} and report its disagreement with the CPU path.

\section{Method}
\label{sec:method}
\begin{figure*}[t]
\centering
\fig{6.2cm}{Figure 1 placeholder: overview of \ours{}. Frozen CLIP ViT-L/14 patch tokens from layers 6/12/18/24 pass through visual adapters and per-layer projections into the text space. Six text anchors (intact, crack, corrosion, spalling, efflorescence, generic defect). LRM mixes patch-level and residual-weighted regional similarity per anchor. AMR reads four independent sigmoids from type-minus-intact margins. RRG computes the residual to one intact reference, yields the defect mass $\nu$, and gates the class-agnostic channel and the image head. Outputs: four type maps, a class-agnostic map, the binary map and the image score.}
\caption{Overview of \ours{}. The backbone is frozen; the 6.6 M trainable parameters sit in the adapters, projections, residual heads and anchor offsets.}
\label{fig:overview}
\end{figure*}

\subsection{Setting and notation}
A query image $x$ is resized to $518\times518$ and encoded by a frozen CLIP ViT-L/14. From layers $l\in\{6,12,18,24\}$ we take the $N=37\times37$ patch tokens $F_l\in\mathbb{R}^{N\times d}$, $d=1024$. A two-branch residual MLP adapter and a per-layer linear projection $P_l:\mathbb{R}^d\to\mathbb{R}^{768}$ map the tokens into the text space, as in AdaptCLIP~\cite{adaptclip}. Supervision consists of a binary mask $Y_\cup$ and four possibly overlapping type masks $Y_c$, $c\in\{1,\dots,4\}$, for the annotated training tiles, and $Y_\cup=0$ for intact tiles. At inference the head receives $x$ and one arbitrary intact image $n$.

\subsection{Text anchors}
Six anchors $t_a\in\mathbb{R}^{768}$, $a\in\{0,\dots,5\}$, represent intact concrete, crack, corrosion, spalling, efflorescence and generic defect. Each anchor is the CLIP text embedding of four hand-written prompt templates, averaged and normalised, plus a bounded learned offset:
\begin{equation}
t_a=\operatorname{norm}\!\left(t_a^{\mathrm{hand}}+\varepsilon\,\Delta_a/\max(1,\|\Delta_a\|)\right),\quad \varepsilon=0.5,
\end{equation}
which keeps the cosine between $t_a$ and its hand-written prompt above $0.866$. The anchors are the only text-side parameters; the text encoder is used once, offline.

\subsection{Local-Regional Similarity Mixture}
The patch-level similarity of token $i$ to anchor $a$ is the cosine in the text space, averaged over the four layers and divided by a temperature $\tau=0.07$:
\begin{equation}
S^{\mathrm{loc}}_{i,a}=\frac{1}{4}\sum_l \cos\!\left(P_l F_{l,i},\,t_a\right)/\tau .
\end{equation}
A regional similarity pools $S^{\mathrm{loc}}$ over a $5\times5$ window with weights that combine spatial proximity and the residual map $\bar A$ of Section~\ref{sec:rrg} downsampled to the grid, so that tokens inside a suspected defect borrow evidence from one another and intact background does not take part:
\begin{equation}
S^{\mathrm{reg}}_{i,a}=\frac{\sum_{j\in\mathcal{N}_i} w_{ij}\,S^{\mathrm{loc}}_{j,a}}{\sum_{j\in\mathcal{N}_i} w_{ij}},\qquad w_{ij}=\bar A_j\exp\!\left(-\tfrac{\|p_i-p_j\|^2}{2\sigma^2}\right),
\end{equation}
with $\sigma=1.5$ cells. Each anchor owns a learned gate $\gamma_a=\operatorname{sigmoid}(g_a)$ that mixes the two:
\begin{equation}
\tilde S_{i,a}=\gamma_a S^{\mathrm{loc}}_{i,a}+(1-\gamma_a)S^{\mathrm{reg}}_{i,a}.
\end{equation}
After training, the crack and intact anchors settle at $\gamma\approx0.59$ while spalling and efflorescence settle at $0.22$ and $0.18$: the thin class keeps its local evidence and the area classes use context. Section~\ref{sec:ablation} shows what fixing $\gamma$ to 0 or 1 costs.

\subsection{Anchor Margin Readout}
The probability that token $i$ carries type $c$ is read from the margin between the type anchor and the intact anchor,
\begin{equation}
p_{i,c}=\operatorname{sigmoid}\!\left(s_c(\tilde S_{i,c}-\tilde S_{i,0})+b_c\right),\quad c=1,\dots,4,
\end{equation}
with a learned scale $s_c$ and bias $b_c$ per type. The four readouts are independent, so a token inside a spalled region may be positive for spalling and corrosion at once. A six-way softmax $\pi_i=\operatorname{softmax}(\tilde S_i)$ is kept alongside for the defect mass $\nu_i=\sum_{a=1}^{5}\pi_{i,a}$, the probability that the token is damaged at all, which the gate below uses.

\subsection{Reference-Residual Gate}
\label{sec:rrg}
For each layer and token, the residual $r_{l,i}=|q_{l,i}-n^*_{l,i}|$ is the element-wise distance between the adapted query token and its nearest adapted token in the intact reference $n$. The residual heads of AdaptCLIP take $[q_{l,i};r_{l,i}]$ from the four layers and produce a residual map $A_{\mathrm{res}}\in[0,1]^{518\times518}$; its downsampled version $\bar A$ enters the regional weights above. The defect mass gates two outputs. The class-agnostic defect map is $A_{\mathrm{cls}}=A_{\mathrm{res}}\odot\operatorname{up}(\nu)$, and the image head pools the residual features with per-token weights $0.5+0.5\nu_i$ before a two-layer MLP. The binary map is $A_{\mathrm{bin}}=\max\!\left(A_{\mathrm{cls}},\max_c \operatorname{up}(p_c)\operatorname{up}(\nu)\right)$, and the image score averages the MLP probability with the mean of the top 1\% of $A_{\mathrm{bin}}$. The residual therefore does not score anomalies; it tells the type heads where intact texture should be ignored, and it catches damage that none of the four anchors claims. Retrieval of the reference is unnecessary: one random intact image and the four most similar images give the same accuracy (Section~\ref{sec:ablation}).

\subsection{Training}
Patch-level targets mark a cell as damaged when any type mask covers at least a quarter of it. The loss is
\begin{equation}
\mathcal{L}=\mathcal{L}_{\mathrm{res}}+\mathcal{L}_{\mathrm{img}}+\mathcal{L}_{\mathrm{grp}}+0.5\,\mathcal{L}_{\mathrm{cls}}+\mathcal{L}_{\mathrm{type}},
\end{equation}
where $\mathcal{L}_{\mathrm{res}}$ and $\mathcal{L}_{\mathrm{cls}}$ are focal plus Dice losses of $A_{\mathrm{res}}$ and $A_{\mathrm{cls}}$ against $Y_\cup$, $\mathcal{L}_{\mathrm{img}}$ is the image-level cross-entropy, and $\mathcal{L}_{\mathrm{type}}$ averages focal plus Dice of the gated type maps against $Y_c$. The grouped cross-entropy $\mathcal{L}_{\mathrm{grp}}$ acts on $\pi$: a damaged cell pays $-\log\nu_i$ plus a cross-entropy inside the defect group whose soft target is shared equally among the types present, and a background cell pays $-\log\pi_{i,0}$. With probability $0.25$ a type is dropped from the target of a batch and its cells are relabelled as generic defect, which trains the fifth anchor to represent damage of unknown type. Visual-side parameters and anchor offsets are updated in alternating epochs for 20 epochs with AdamW at learning rate $10^{-3}$ and an effective batch of 8. Backbone features are cached once, so training takes about 1.5 hours on one 16 GB GPU.

\section{Experiments}
\label{sec:exp}
\subsection{Data, protocol and compared methods}
\textbf{CODEBRIM tiles.} We cut $512\times512$ tiles from CODEBRIM~\cite{codebrim} images and annotated pixel masks for crack, corrosion, spalling and efflorescence on 917 damaged tiles, keeping overlapping regions as members of every type they belong to. The training split holds these 917 tiles and 1{,}944 intact tiles. The validation split (400 tiles, 123 damaged) served every design choice and the audit of Section~\ref{sec:protocol}. The test split (809 tiles, 254 damaged, 555 intact) was never used for any decision; all main tables report it. The two splits differ in defect size: the median damaged tile carries 6{,}297 defect pixels in validation and 44{,}114 in test, where spalling and efflorescence dominate.

\textbf{External sets.} S2DS~\cite{s2ds} test (447 tiles) and a strict four-type subset of dacl10k~\cite{dacl10k} (262 tiles, 207 damaged) are evaluated with the CODEBRIM-trained models and a CODEBRIM intact tile as reference. Per-type ground truth is decoded from their flattened masks, so overlapping damage counts for one class there.

\textbf{Metrics.} Image-level AUROC (I-AUROC); pixel-level AUROC, AP, F1 at the best threshold and AUPRO with a 0.3 FPR limit on the union mask; per-type AP and best-threshold IoU on each type mask, averaged into mAP and mIoU; and the fraction of intact-tile area flagged at the threshold that reaches a given pixel recall on damaged tiles, which measures false alarms on intact surface at matched sensitivity. AUPRO is computed with the GPU implementation of anomalib~\cite{anomalib}; its CPU path gives 36.7 instead of 41.5 on identical S2DS predictions, so all rows use one implementation.

\textbf{Compared methods.} All methods are trained on the same training tiles with their official code: AdaptCLIP~\cite{adaptclip} (1.8 M trainable parameters), AnomalyCLIP~\cite{anomalyclip} (5.6 M), Bayes-PFL~\cite{bayespfl} (27.4 M), MultiADS~\cite{multiads} (3.1 M; the official learning rate of $10^{-3}$ collapsed to all-intact predictions on these tiles and $10^{-4}$ was used), SegFormer-B2~\cite{segformer} (25 M, multi-label sigmoid output on the four type masks, 40 epochs) and WPFormer~\cite{wpformer} (26.1 M, binary union mask, 60 epochs, checkpoint selected on validation). SegFormer-B2 and MultiADS have per-type outputs; the others produce a binary map. \ours{} is reported as the mean of three seeds with standard deviation; the other methods are single runs.

\subsection{In-domain results}
\begin{table*}[t]
\centering
\caption{Binary detection and localisation on the CODEBRIM test split (809 tiles). FA@0.8/0.9: fraction of intact-tile area (\%) flagged at the threshold giving 80\% / 90\% pixel recall on damaged tiles. \ours{}: mean $\pm$ std over three seeds; other methods single runs.}
\label{tab:binary}
\small
\setlength{\tabcolsep}{6pt}
\begin{tabular}{lrcccccc}
\toprule
Method & Param. & I-AUROC & P-AUROC & P-AP & P-AUPRO & FA@0.8 & FA@0.9 \\
\midrule
AnomalyCLIP & 5.6 M & 82.3 & 92.0 & 57.0 & 71.0 & 11.3 & 18.4 \\
Bayes-PFL & 27.4 M & 91.6 & 93.2 & 62.7 & 75.6 & 8.6 & 16.1 \\
MultiADS & 3.1 M & 72.3 & 93.7 & 69.8 & 78.8 & 6.5 & 12.0 \\
AdaptCLIP & 1.8 M & \best{92.7} & 95.5 & 73.3 & \best{80.4} & 4.6 & 9.5 \\
SegFormer-B2 & 25 M & 90.2 & 94.7 & 76.4 & 68.6 & \best{2.4} & 10.0 \\
WPFormer & 26.1 M & 88.8 & 91.0 & 70.8 & 78.0 & 11.1 & 27.7 \\
\ours{} & 6.6 M & 92.0$\pm$0.7 & \best{96.3}$\pm$0.5 & \best{76.8}$\pm$3.1 & 73.5$\pm$3.6 & 2.5 & \best{4.4} \\
\bottomrule
\end{tabular}
\end{table*}

Table~\ref{tab:binary} reports binary detection and localisation on the test split. \ours{} reaches a pixel AP of 76.8, level with SegFormer-B2 at 76.4 and above AdaptCLIP by 3.5, MultiADS by 7.0, WPFormer by 6.0 and Bayes-PFL by 14.1 points. Its pixel AUROC of 96.3 is the highest in the table. Image-level AUROC is 92.0 against 92.7 for AdaptCLIP, within one standard deviation. The false-alarm columns separate the methods most clearly. At 80\% pixel recall \ours{} and SegFormer-B2 flag 2.5\% and 2.4\% of intact surface while the CLIP detectors flag 4.6\% to 11.3\%; at 90\% recall SegFormer-B2 rises to 10.0\% and WPFormer to 27.7\%, whereas \ours{} stays at 4.4\% over all three seeds (4.3 to 4.6). Reaching the last ten points of recall costs the supervised networks a four-fold increase in false-alarm area and \ours{} a 1.8-fold increase. P-AUPRO orders the methods differently: AdaptCLIP at 80.4 and MultiADS at 78.8 are above \ours{} at 73.5. P-AUPRO averages the recall of every connected component and so weights each small crack fragment as much as a large spalled region; the false-alarm area at matched recall measures the opposite side of the trade-off, and \ours{} sits on the low-false-alarm side.

\begin{table}[t]
\centering
\caption{Per-type segmentation on the CODEBRIM test split for the three methods with type outputs. AP per type, mAP and best-threshold mIoU in \%.}
\label{tab:pertype}
\small
\setlength{\tabcolsep}{3.5pt}
\begin{tabular}{lcccccc}
\toprule
Method & Crack & Corrosion & Spalling & Efflor. & mAP & mIoU \\
\midrule
MultiADS & 15.2 & 79.9 & 71.0 & 58.7 & 56.2 & 40.2 \\
SegFormer-B2 & \best{24.1} & \best{81.8} & \best{76.1} & 58.5 & \best{60.1} & \best{50.4} \\
\ours{} & 10.3 & 80.8 & 76.0 & \best{70.5} & 59.4 & 40.7 \\
\bottomrule
\end{tabular}
\end{table}

Table~\ref{tab:pertype} gives the per-type results. \ours{} and SegFormer-B2 are within one point on corrosion and spalling; \ours{} leads on efflorescence by 12 points, and SegFormer-B2 leads on cracks by 14 points, which decides the mAP of 59.4 against 60.1. MultiADS is 3.2 mAP below \ours{} with a lower spalling and efflorescence AP. Cracks are the weak type for every method on this split: they cover 1.8\% of the defect pixels, and the thin structure is below the $37\times37$ resolution of the CLIP grid. The mIoU gap of ten points to SegFormer-B2 comes from the same source, since IoU at a fixed threshold penalises the blurred boundaries of a 14-pixel grid more than AP does.

\subsection{Ablation}
\label{sec:ablation}
\begin{table}[t]
\centering
\caption{Ablation of \ours{} on the CODEBRIM test split (seed 10). Val crack AP in the last column is the validation-split crack AP, where cracks are frequent enough to separate the variants. Differences below the seed standard deviation of Table~\ref{tab:binary} (3.1 P-AP, 3.6 P-AUPRO) are not interpreted.}
\label{tab:ablation}
\scriptsize
\setlength{\tabcolsep}{2.6pt}
\resizebox{\linewidth}{!}{%
\begin{tabular}{lcccccc}
\toprule
Variant & I-AUROC & P-AP & P-AUPRO & mAP & FA@0.8 & Val crack \\
\midrule
Full \ours{} & 91.4 & 76.3 & 70.1 & 58.7 & 2.4 & 35.7 \\
LRM: $\gamma=0$ (regional only) & 90.0 & 74.4 & 69.9 & 57.3 & 2.5 & 8.2 \\
LRM: $\gamma=1$ (local only) & 90.6 & 75.0 & 72.0 & 59.4 & 2.3 & 37.6 \\
AMR $\to$ 13-way softmax (v0) & 88.6 & 71.4 & 65.6 & 48.6 & 6.0 & 6.8 \\
RRG removed (no residual) & 89.7 & 73.2 & 71.2 & 57.2 & 2.9 & 34.3 \\
RRG without $\nu$ gating & 91.8 & 76.2 & 71.8 & 59.7 & 3.0 & 35.5 \\
12 anchors (6 nuisance anchors) & 91.0 & 72.1 & 73.5 & 56.9 & 2.8 & 40.5 \\
\bottomrule
\end{tabular}}
\end{table}

Table~\ref{tab:ablation} isolates the three components. The readout matters most. Replacing the independent margins by the 13-way softmax of our first design costs 4.9 pixel AP and 10.1 mAP, raises the false-alarm area from 2.4\% to 6.0\%, and leaves crack AP at 6.8 on the validation split, because the softmax hands a thin crack to the intact anchor whenever the regional evidence is weak. The local branch of LRM is what keeps cracks: fixing $\gamma=0$ drops validation crack AP from 35.7 to 8.2 and test crack AP from 8.3 to 5.2 while the area classes move by less than two points; fixing $\gamma=1$ leaves the typed outputs intact but removes the regional evidence that the gated outputs rely on, so the gated mAP falls from 58.4 to 49.0. Removing the residual branch costs 3.1 pixel AP on the test split, within seed variance, and 9.7 on the validation split, where small defects on intact texture are the hard case and the false-alarm area on cleaned intact tiles grows from 0.12\% to 0.93\%. Removing only the $\nu$ gating changes no test metric beyond seed variance; the gate stays in the model for the class-agnostic channel and the unknown-type use, not for a measured accuracy gain. Six additional anchors naming nuisances such as shadow, water stain or repair patch receive 0.1\% of the arg-max assignments; over three seeds on the validation split the 12-anchor and 6-anchor models have the same mean pixel AP and AUPRO, and the single 12-anchor run on the test split is 4.2 pixel AP lower and 3.4 P-AUPRO higher, so the nuisance anchors were dropped. Using one random intact tile instead of the four retrieved most similar tiles changes pixel AP by 0.2 on the validation split, so inference needs no reference library.

\subsection{Cross-dataset results}
\begin{table*}[t]
\centering
\caption{Cross-dataset evaluation of the CODEBRIM-trained models on S2DS test (447 tiles) and the dacl10k four-type subset (262 tiles), without target-domain data. mAP only for methods with type outputs.}
\label{tab:cross}
\small
\setlength{\tabcolsep}{6pt}
\begin{tabular}{lcccc|cccc}
\toprule
& \multicolumn{4}{c|}{S2DS} & \multicolumn{4}{c}{dacl10k} \\
Method & I-AUROC & P-AP & P-AUPRO & mAP & I-AUROC & P-AP & P-AUPRO & mAP \\
\midrule
AnomalyCLIP & 71.9 & 53.0 & 55.1 & -- & 53.4 & 26.1 & 52.1 & -- \\
Bayes-PFL & \best{85.1} & 56.3 & 50.9 & -- & \best{69.9} & 31.5 & 50.9 & -- \\
MultiADS & 76.2 & \best{71.5} & 61.4 & \best{52.8} & 67.4 & 36.2 & 49.6 & \best{18.5} \\
AdaptCLIP & 81.3 & 62.9 & 63.7 & -- & 67.3 & \best{40.6} & \best{54.7} & -- \\
SegFormer-B2 & 74.3 & 63.0 & 45.7 & 32.7 & 56.2 & 22.0 & 36.6 & 8.6 \\
WPFormer & 84.6 & 64.4 & \best{64.2} & -- & 67.4 & 28.4 & 37.6 & -- \\
\ours{} & 78.1 & 70.9 & 41.5 & 46.6 & 69.8 & 32.8 & 40.1 & 16.1 \\
\bottomrule
\end{tabular}
\end{table*}

Table~\ref{tab:cross} moves the same checkpoints to S2DS and dacl10k. The two type-aware frozen-CLIP heads lead in pixel AP on both sets: \ours{} and MultiADS reach 70.9 and 71.5 on S2DS against 64.4 for WPFormer and 63.0 for SegFormer-B2, and 32.8 and 36.2 on dacl10k against 28.4 and 22.0. In per-type segmentation \ours{} is 14 mAP above SegFormer-B2 on S2DS and 7.5 above on dacl10k, and MultiADS is above \ours{} by 6.2 and 2.4. The ordering of Table~\ref{tab:binary} therefore inverts: SegFormer-B2, level with \ours{} in-domain, loses 13 points of mAP relative to it on S2DS, while MultiADS, 7 points behind in-domain, moves ahead. Fine-tuning \ours{} for ten more epochs shifts its external numbers by two to five points in inconsistent directions, so differences of that size between frozen-CLIP methods are within training noise, and the gap to the supervised networks is not. P-AUPRO and image-level AUROC do not follow this pattern: WPFormer has the highest S2DS P-AUPRO at 64.2 and Bayes-PFL the highest I-AUROC on both sets, and \ours{} has the lowest S2DS P-AUPRO at 41.5, which reflects its low recall on the many small crack components of S2DS. The components that bring in-domain accuracy, the residual branch and the supervised type heads, are fitted to the appearance of CODEBRIM; what transfers is the frozen backbone and its alignment with the text anchors, which \ours{} and MultiADS share with a smaller supervised footprint in the latter.

\subsection{Evaluation protocol}
\label{sec:protocol}
Three measurement choices change the conclusions above. Reviewing the 277 intact validation tiles ranked by model score found 21 tiles with unlabelled damage; removing them (Cleaned-Val) lowers the false-alarm area of \ours{} at 80\% recall from 1.75\% to 0.12\% and raises its validation pixel AP from 81.8 to 89.9 without changing the ranking of methods. The CPU and GPU implementations of AUPRO in anomalib disagree by 4.8 points on identical S2DS predictions because their connected-component labelling differs. The defect-size distribution of the split has the largest effect: the same checkpoints give a 20-point pixel-AP margin of \ours{} over AdaptCLIP on the small-defect validation split and 3.5 points on the large-defect test split, while the margin over SegFormer-B2 moves from $-6.3$ to $+0.4$. Comparisons between methods in this field are only meaningful on a stated split with a stated defect-size distribution and a stated AUPRO implementation.

\subsection{Crack scale}
Resolution explains part of the crack gap. Sliding-window inference at twice the resolution raises validation crack AP from 35.7 to 42.3 but lowers efflorescence AP by 25 points, since the regional context that efflorescence needs falls outside the window. A quarter-resolution convolutional detail branch trained as a third scale reaches 42.4 on cracks when trained alone and 37.8 when trained jointly, and neither variant transfers to S2DS or dacl10k. The thin class needs resolution that the $37\times37$ grid does not provide, and recovering it in-domain does not change the cross-dataset picture.

\subsection{Qualitative results}
\begin{figure}[t]
\centering
\fig{4.8cm}{Figure 2 placeholder: qualitative results. Rows: CODEBRIM test, S2DS, dacl10k. Columns: image, per-type ground truth, \ours{} type maps (crack / corrosion / spalling / efflorescence overlaid with overlap shown), class-agnostic map $A_{\mathrm{cls}}$, SegFormer-B2, MultiADS.}
\caption{Qualitative results on the three test sets. Overlapping types are shown as mixed colours; the class-agnostic map responds to damage that no type anchor claims.}
\label{fig:qual}
\end{figure}
Figure~\ref{fig:qual} shows type maps on the three test sets, including spalled regions with corrosion inside them that the independent readouts mark as both, and responses of the class-agnostic channel on damage outside the four types.

\section{Conclusion}
\ours{} uses the text anchors of a frozen CLIP as a class interface for a light supervised head: independent readouts on anchor margins give overlapping per-type masks, a per-anchor local-regional mixture keeps thin cracks visible, and a residual gate from one intact image suppresses responses on intact texture. On a held-out CODEBRIM test split the head matches a fully trained SegFormer-B2 in binary pixel AP, exceeds every CLIP-based detector trained on the same tiles, and flags less than half the intact-surface area of any other method at 90\% recall. Across datasets, type-aware frozen-CLIP heads exceed both supervised networks and the simpler MultiADS exceeds \ours{}. In-domain accuracy of a supervised head is accuracy on the training domain; transfer comes from the frozen backbone and its alignment with text. The per-type masks and the evaluation protocol are available on request.

\bibliographystyle{IEEEtran}
\bibliography{refs}

\end{document}
